% translation-source: docs/user/latex/chapters/03-informations-selection.tex; sha256: d6262e09b4a7770461e960518bb3c510f2e2d54c066a2e8584fc3ce8f21c3dd8
\chapter{Information, rating and selection}
\section{Read image information}
The left panel describes the open photo: name, available date, storage,
dimensions, camera and other known metadata. Missing information must not be
interpreted as a measured zero. Available shooting information below the
histogram helps you read ISO, focal length and exposure time.

\section{Rate and flag a photo}
Click one of the five stars to assign a rating from 1 to 5. Clicking the selected
rating again resets it to zero. Click to toggle the flag. Choose a red,
yellow, green, blue or purple label; clicking the same color removes it.
These organizational choices do not recolor pixels and are not mixer bands.

No numeric rating or flag shortcut is currently declared. Use the Information
panel controls.

\section{History and snapshots}
Under the labels, the left panel lists the steps of the open photo, the most
recent at the top. A drag of a slider or wheel, a moved curve point, a typed
value or a profile choice each count as one step. The last hundred steps are
saved in the photo's sidecar and come back when it is reopened.

\emph{Undo} and \emph{Redo} are in this section and in the \emph{Edit} menu.
Clicking a step returns to it; more recent steps appear dimmed and stay
available until the next edit. Undoing a profile change develops the photo
again.

To remove an edit from the middle of the history, right-click its step and
choose \emph{Delete this step}, or press \touche{Delete} when the list has the
focus: the current step is then removed. Later steps are replayed without this
edit and keep what they changed themselves. The opening step cannot be removed.

A snapshot keeps the current state under a name. Type the name, then choose
\emph{Create}; clicking the snapshot restores it as a new step.

\attention{Once saved by this version, a sidecar can no longer be read by an
older version of IPR. A sidecar written by a more recent version opens as
read-only.}

\section{Copy and paste settings}
On a ready image, choose \emph{Edit > Copy settings} or
\touche{Ctrl + Shift + C}. Open the destination photo, wait for processing,
then choose \emph{Paste settings} or \touche{Ctrl + Shift + V}.
These commands are also available by right-clicking the open image.
The copy is a snapshot: subsequent source edits do not alter it. It remains
available during the session, even after switching folders.

Pasting replaces all edits, including curves, mixer and grading. It preserves
the destination profile, calibration, optics, automatic white balance and
rating. If the source uses as-shot white balance, the destination uses its
own as-shot balance; custom temperature and tint are copied as absolute
values. Two cameras will not necessarily produce the same color.
There is no selection of individual groups to paste yet; a paste can be undone
like any other history step.
\touche{Ctrl + C} and \touche{Ctrl + V} remain available in text fields.

\section{Copy edits to multiple photos}
\begin{enumerate}
\item Open the photo whose settings you want to copy.
\item Check destination thumbnails or use \touche{Ctrl + click}; do not
confuse a checked box with the open image. The counter shows selected files.
\item Right-click and choose \emph{Apply the open image's settings to the selection},
also in the \emph{Edit} menu. Check the result message.
\item Uncheck files or use the clear-selection button.
\end{enumerate}
This writes edits into destination sidecars. It is not limited to the currently
expanded group. Destination RAW profiles and automatic white balance are not
blindly copied between cameras; destination organizational metadata is preserved.

\attention{This is batch application of edits, not the multiple-export
command. \touche{Ctrl + click} only toggles the targeted thumbnail without
opening it; an ordinary click opens it without clearing the selection.
Shift-click range selection is not implemented.}

\section{Filmstrip badges}
The open thumbnail's badge may show rating, flag and label. No badge on another
thumbnail does not mean that it has no rating or edits: those data are not yet
exposed for every model row.

\section{Right-click a thumbnail}
The menu targets the thumbnail under the pointer. For a photo other than the
open one, first choose \emph{Open this image}: right-clicking alone neither
loads it nor changes multiple selection. Editing actions remain disabled
until that target is open, avoiding an accidental reset or export of the
previous photo.

\emph{Duplicate file and settings}, associated with the reserved
\touche{Ctrl + D} shortcut, is intended to make a physical copy and a consistent
sidecar, not a developed TIFF or virtual copy. This command still awaits its
engine interface and is disabled in the version described.
